\documentclass[aps,prl,twocolumn,showpacs,superscriptaddress]{revtex4-2}
\usepackage{amsmath,amssymb,booktabs,xcolor,hyperref,amsthm}
\newtheorem{theorem}{Theorem}
\newtheorem{definition}{Definition}
\newtheorem{corollary}{Corollary}

\begin{document}

\title{The Two Fundamental Theorems of BCT Chemistry:\\
Exact Proofs of the Pauli Exclusion Principle and Noble Gas Inertness
from Hopf Charge Quantisation}

\author{Michel Robert Cabri\'e}
\affiliation{Independent Researcher, Victoria, Australia}
\email{ZeroFreeParameters@gmail.com}
\date{20 March 2026}

\begin{abstract}
Letter~78 introduced atomic structure as Hopf topology and stated the 
Pauli exclusion principle and noble gas inertness as consequences of BCT 
vacuum geometry. This Letter provides the exact proofs. Both theorems 
derive from a single principle: Hopf charge is integer-valued 
($H\in\mathbb{Z}$, from $\pi_3(S^2)=\mathbb{Z}$) and the $n$-th OHC 
Bessel resonance mode has capacity $2n^2$.
\textbf{Pauli Exclusion Theorem:} The $n$-th OHC mode supports $|H|\leq n$.
Two $H=+1$ Hopfions in the $n=1$ mode give $|H|=2>1$, topologically 
forbidden. $H=+1$ and $H=-1$ give $|H|=0\leq1$, allowed. The exclusion 
principle is a theorem of OHC mode capacity, not a postulate.
\textbf{Noble Gas Theorem:} A closed Hopf multiplet --- all $2n^2$ modes 
filled, $H_{\rm net}=0$, no vacant modes --- cannot form chemical bonds 
because bonding requires a vacant mode or unpaired Hopfion, and a closed 
multiplet has neither. Noble gas inertness is topologically exact, not 
merely energetically favourable. Both proofs are rigorous, exact, 
and parameter-free.
\end{abstract}

\maketitle

\section{Introduction}

Letter~78 derived atomic structure from BCT vacuum topology and stated 
the Pauli exclusion principle and noble gas inertness as consequences. 
This Letter provides the exact proofs, showing both are strict theorems 
of the Octet-Hopfion Condensate (OHC) topology established in 
Appendix~JH~\cite{AppJH}.

\section{Mathematical Preliminaries}

\subsection{Hopf Charge Quantisation}

The Hopf fibration $\eta: S^3\to S^2$ classifies maps by 
$\pi_3(S^2)=\mathbb{Z}$. The Hopf invariant $H\in\mathbb{Z}$ is additive:
\begin{equation}
    H(f \circ g) = H(f) + H(g)
\end{equation}
Each electron is identified with the $H=1$ generator. Two electrons 
with the same OHC phase orientation combine to $H=+2$; opposite 
orientations combine to $H=0$.

\subsection{OHC Mode Capacity}

The $n$-th OHC Bessel resonance mode supports Hopf charge 
$|H_{\rm total}|\leq n$, with $n^2$ spatial modes each 
admitting two spin orientations, giving total capacity $2n^2$.

\section{The Pauli Exclusion Theorem}

\begin{theorem}[BCT Pauli Exclusion]
No two Hopfions in the same OHC Bessel resonance mode may have 
the same phase orientation.
\end{theorem}

\begin{proof}
Let two Hopfions occupy the same spatial mode at level $n=1$.

\textit{Case 1 --- Same orientation.} Both carry $H=+1$. Combined 
charge $H_{\rm total} = +1+1 = +2$. Since $|{+2}|=2>n=1$, this 
exceeds mode capacity. Topologically forbidden.

\textit{Case 2 --- Opposite orientation.} Charges $H=+1$ and $H=-1$. 
Combined $H_{\rm total} = 0$. Since $|0|=0\leq1$, this satisfies 
capacity. Topologically allowed.

Therefore at most two Hopfions may coexist in one mode, and they 
must have opposite phase orientations (opposite spin). \qed
\end{proof}

\begin{corollary}
Electron spin is the topological winding direction of the Hopf 
map $S^3\to S^2$: $H=+1$ (spin up) and $H=-1$ (spin down). 
It is quantised because $\pi_3(S^2)=\mathbb{Z}$.
\end{corollary}

\section{The Noble Gas Theorem}

\begin{definition}[Closed Hopf Multiplet]
A closed Hopf multiplet at level $n$ is a configuration of $2n^2$ 
Hopfions that: (a) fills all $n^2$ spatial modes at level $n$; 
(b) occupies both spin orientations in each mode; 
(c) has $H_{\rm net}=0$; (d) has no vacant modes; 
(e) has no unpaired Hopfions.
\end{definition}

\begin{theorem}[BCT Noble Gas]
An atom whose outer shell is a closed Hopf multiplet cannot 
form chemical bonds.
\end{theorem}

\begin{proof}
Chemical bonding requires one of:
\begin{itemize}
\item[(i)] Covalent: sharing a Hopf mode $\to$ requires vacant 
mode or unpaired Hopfion in A's outer shell.
\item[(ii)] Ionic: transferring a Hopfion $\to$ requires unpaired 
Hopfion (donor) or vacant mode (acceptor).
\end{itemize}
Both require condition (d) or (e) to be violated. A closed 
Hopf multiplet satisfies (a)--(e) by definition; in particular 
it has no vacant modes and no unpaired Hopfions. 
Therefore it cannot form bonds. \qed

\textit{Remark.} This is topological impossibility, not energetic 
unfavourability. Energetic barriers can be overcome; topological 
impossibilities cannot.
\end{proof}

Noble gas ionisation energies (24.6, 21.6, 15.8, 14.0, 12.1, 
10.7~eV for He through Rn) are $600$--$1000\,k_BT$ at room 
temperature, consistent with extreme topological stability.

\section{The Common Root}

Both theorems derive from one principle:
\begin{equation}
    H \in \mathbb{Z} \quad (\pi_3(S^2)=\mathbb{Z})
    \quad + \quad \text{capacity } 2n^2
\end{equation}

Pauli exclusion: $|H|\leq n$ per mode $\to$ at most 2 electrons 
per spatial mode, opposite spin.

Noble gas: closed multiplet has no vacant modes, no unpaired 
Hopfions $\to$ bonding topologically impossible.

All of chemistry rests on two integers and one geometry.

\section{Comparison with Standard QM}

\begin{center}
\begin{tabular}{lll}
\toprule
Question & Standard QM & BCT \\
\midrule
Pauli exclusion & Postulated axiom & Theorem 1 (exact) \\
Spin & Postulated/Dirac & Hopf winding $H=\pm1$ \\
Shell capacity $2n^2$ & Derived, unexplained & $n^2$ Hopf modes $\times2$ \\
Noble gas inertness & Energetic & Topological (Theorem 2) \\
Why bonding occurs & Energy lowering & Open Hopf channels \\
\bottomrule
\end{tabular}
\end{center}

\section{Conclusion}

The Pauli exclusion principle and noble gas inertness are exact 
theorems of BCT vacuum topology. The proofs are rigorous, exact, 
and parameter-free. Spin is the Hopf winding direction. Shell 
capacity is Hopf mode counting. Noble gas inertness is topological 
impossibility. All of chemistry rests on $H\in\mathbb{Z}$ and 
capacity $2n^2$.

\begin{acknowledgments}
This Letter provides exact proofs of results stated in Letter~78 
(20~March~2026). The proofs emerged from recognising that 
$H\in\mathbb{Z}$ and mode capacity $2n^2$ are sufficient to derive 
all fundamental rules of chemistry without additional postulates.
\end{acknowledgments}

\begin{thebibliography}{99}
\bibitem{AppJH} M.~R.~Cabri\'e, \textit{Appendix JH: The Hopfion 
Interior and OHC}, Zenodo (2026), doi:10.5281/zenodo.18975018
\bibitem{L78} M.~R.~Cabri\'e, \textit{BCT Letter 78: Atomic Structure 
as Hopfion Topology}, Zenodo (2026).
\bibitem{L80} M.~R.~Cabri\'e, \textit{BCT Letter 80: Molecular Geometry 
from Hopf Topology}, Zenodo (2026).
\end{thebibliography}

\end{document}
